% !PRINTED 62
\weberpage{62}
%<<CONT
und nur wenn $\lambda = 2$ ist, ist $c_0 = 2$ und nicht $= 1$ zu setzen.

Wenn wir dann mit $\varepsilon_{-1}, \varepsilon_0, \varepsilon_1 \ldots \varepsilon_\mu$ ein System primitiver
Einheitswurzeln der Grade $c_{-1}, c_0, c_1 \ldots c_\mu$ bezeichnen, und mit
$\beta_{-1}, \beta_0, \beta_1 \ldots \beta_\mu$ die Indices einer Zahl $b$, so erhalten wir die
Charaktere der Gruppe $\mathfrak{N}$ in der Form

\weq{7}{\chi_b(n) = \varepsilon_{-1}^{\beta_{-1}\nu_{-1}} \varepsilon_0^{\beta_0\nu_0} \varepsilon_1^{\beta_1\nu_1} \ldots \varepsilon_\mu^{\beta_\mu\nu_\mu}.}

Darin ist $\varepsilon_{-1} = 1$, wenn $\lambda = 0, 1, 2$ ist; in den anderen
F\"allen ist $\varepsilon_{-1} = -1$, und $\varepsilon_1, \ldots \varepsilon_\mu$ sind primitive Einheitswurzeln
 der Grade
\[
c_1 = \varphi(q_1^{\varkappa_1}), \ldots c_\mu = \varphi(q_\mu^{\varkappa_\mu}).
\]

\begin{center}
\rule{3cm}{0.4pt}
\end{center}
